% !TEX program = xelatex
\documentclass[11pt,a4paper]{article}
\usepackage{fix-cm}
\usepackage[top=20mm,bottom=20mm,left=22mm,right=22mm,headheight=17pt,headsep=7mm]{geometry}
\usepackage{graphicx,amsmath,amssymb,booktabs,array,tabularx,xcolor,float}
\usepackage{tikz}
\usetikzlibrary{arrows.meta,positioning,calc}
\usepackage{caption,fancyhdr,titlesec}
\usepackage[unicode,hidelinks]{hyperref}
\usepackage{xepersian}
\providecommand{\UseMathForPositioningText}{}
\settextfont[Path=fonts/,BoldFont=Sahel.ttf,BoldFeatures={FakeBold=1.5}]{Sahel.ttf}
\setlatintextfont{TeX Gyre Termes}
\setdigitfont[Path=fonts/,BoldFont=Sahel.ttf,BoldFeatures={FakeBold=1.5},ItalicFont=Sahel.ttf,BoldItalicFont=Sahel.ttf]{Sahel.ttf}
\definecolor{navy}{HTML}{285F87}
\definecolor{amber}{HTML}{D18A31}
\definecolor{ink}{HTML}{263542}
\definecolor{muted}{HTML}{667582}
\definecolor{pale}{HTML}{EEF4F8}
\hypersetup{pdftitle={طبقه‌بندی بار ضربان زودرس بطنی با ویژگی‌های بای‌اسپکتروم},pdfauthor={},pdfsubject={گزارش کار تحلیل سری ضربان قلب}}
\titleformat{\section}{\large\bfseries\color{navy}}{\thesection.}{0.45em}{}
\titlespacing*{\section}{0pt}{7pt}{4pt}
\captionsetup{font=small,labelfont=bf,justification=centering,skip=4pt}
\setlength{\parindent}{0pt}
\setlength{\parskip}{4pt}
\setlength{\emergencystretch}{2em}
\renewcommand{\arraystretch}{1.18}
\setlength{\tabcolsep}{5pt}
\pagestyle{plain}
\fancyhf{}
\fancyhead[R]{\small\color{muted}گزارش کار تحلیل سری ضربان قلب}
\fancyhead[L]{\small\color{muted}\lr{HOS / MIT-BIH}}
\fancyfoot[C]{\small\thepage}
\renewcommand{\headrulewidth}{0.3pt}
\renewcommand{\footrulewidth}{0pt}
\newcommand{\fnum}[1]{\lr{\persianfont #1}}
\newcommand{\interval}[2]{\lr{[#1, #2]}}
\newcommand{\figwide}[2]{\includegraphics[width=\linewidth,height=#2,keepaspectratio]{images/#1}}
\begin{document}
\fontsize{11.2}{15.1}\selectfont
\thispagestyle{plain}
\begin{center}
{\LARGE\bfseries\color{navy}طبقه‌بندی بار ضربان زودرس بطنی\par}
\vspace{2pt}
{\LARGE\bfseries\color{navy}با ویژگی‌های بای‌اسپکتروم\par}\vspace{6pt}
{\normalsize مقایسهٔ ویژگی‌های مرتبهٔ بالاتر با شاخص‌های سادهٔ سری ضربان قلب}
\end{center}
\vspace{-2mm}
\section*{چکیده}
\begingroup\small
در این کار بررسی شد که آیا ویژگی‌های بای‌اسپکترومِ سری ضربان قلب می‌توانند قطعات دارای بار کم و زیاد ضربان زودرس بطنی \lr{(PVC)} را در افراد خارج از آموزش تفکیک کنند. داده‌ها از پایگاه عمومی \lr{MIT-BIH Arrhythmia} انتخاب شدند. از میان ۹۶ قطعهٔ ۱۵ دقیقه‌ای، ۳۸ قطعه از ۲۲ فرد شرایط ورود را داشتند؛ ۲۶ قطعه کم‌بار و ۱۲ قطعه پربار بودند. بار حداکثر \fnum{۱٪} و حداقل \fnum{۵٪} به‌ترتیب دو کلاس را تعریف کردند. زمان ضربان‌ها از دو مسیر مرجع و آشکارسازی از نوار قلب به دست آمد و پس از تبدیل به سری ضربان با نرخ ۴ هرتز، هفت ویژگی طیفی مرتبهٔ بالاتر \lr{(HOS)} استخراج شد. هشت ویژگی ساده و مجموعهٔ ترکیبی نیز مقایسه شدند. ارزیابی با اعتبارسنجی متقاطع تو‌در‌توی فردمحور و ۲۰۰۰ بازنمونه‌گیری در سطح فرد انجام شد. در مسیر مرجع، دقت متوازنِ ویژگی‌های ساده، \lr{HOS} و مجموعهٔ ترکیبی به‌ترتیب \fnum{۹۵٫۸۳}، \fnum{۸۹٫۷۴} و \fnum{۹۳٫۹۱٪} بود. در مسیر آشکارساز، \lr{HOS} به \fnum{۹۶٫۱۵٪} رسید. با توجه به تعداد محدود افراد و بازه‌های اطمینان، این نتایج برتری پایدار \lr{HOS} را نشان نمی‌دهند؛ با این حال، قابلیت استفاده از این ویژگی‌ها برای مسئلهٔ تعریف‌شده را نشان می‌دهند.
\par\endgroup
\textbf{واژه‌های کلیدی:} ضربان زودرس بطنی، فاصلهٔ بین ضربان‌ها، بای‌اسپکتروم، طبقه‌بندی، ارزیابی فردمحور.
\section{مقدمه و هدف}
ضربان زودرس بطنی می‌تواند توالی زمانی ضربان‌ها را تغییر دهد و الگوهایی مانند فاصلهٔ کوتاه و مکث پس از آن ایجاد کند. بنابراین، اطلاعات مربوط به این پدیده تنها در شکل موج نوار قلب خلاصه نمی‌شود و در سری فاصلهٔ بین ضربان‌ها نیز قابل بررسی است. بای‌اسپکتروم، ساختار طیفی مرتبهٔ سوم و رابطهٔ فازی میان مؤلفه‌های بسامدی را توصیف می‌کند. روش مبتنی بر این ویژگی‌ها در مطالعهٔ \lr{Chua} و همکاران برای گروه‌های قلبی مختلف بررسی شده است \cite{chua}.

پرسش این آزمایش آن بود که آیا هفت ویژگی بای‌اسپکتروم، قطعات کم‌بار و پربار \lr{PVC} را در افراد جدید از هم جدا می‌کنند و آیا نسبت به ویژگی‌های سادهٔ سری ضربان، اطلاعات مکمل دارند. واحد تصمیم‌گیری یک قطعهٔ ۱۵ دقیقه‌ای بود. این هدف با تشخیص نوع تک‌تک ضربان‌ها یا بازسازی گروه‌های بالینی مقالهٔ مبنا تفاوت دارد.
\section{داده‌ها و تعریف کلاس‌ها}
پایگاه \lr{MIT-BIH Arrhythmia} شامل ۴۸ رکورد نیم‌ساعته از ۴۷ فرد با نرخ نمونه‌برداری ۳۶۰ هرتز است \cite{mitdb}. از هر رکورد دو قطعهٔ بدون هم‌پوشانی در بازه‌های صفر تا ۹۰۰ و ۹۰۰ تا ۱۸۰۰ ثانیه در نظر گرفته شد. برچسب هر قطعه با تقسیم تعداد ضربان‌های دارای حاشیه‌نویسی \lr{V} بر تعداد کل ضربان‌های واقعی آن قطعه تعیین شد. نشانه‌های ریتم و کیفیت سیگنال در مخرج این نسبت قرار نگرفتند.
\begin{table}[H]
\centering\small
\caption{تعریف عملیاتی کلاس‌ها و جمعیت نهایی آزمایش.}
\begin{tabular}{rccc}
\toprule
کلاس & سهم ضربان‌های \lr{V} & تعداد قطعه & تعداد فرد\\
\midrule
کم‌بار & حداکثر \fnum{۱٪} & ۲۶ & ۱۵\\
پربار & حداقل \fnum{۵٪} & ۱۲ & ۷\\
\bottomrule
\end{tabular}
\label{tab:cohort}
\end{table}
آستانه‌ها برای تعریف این آزمایش انتخاب شدند و معیار تشخیص بالینی نیستند. قطعات بین \fnum{۱٪} و \fnum{۵٪} کنار گذاشته شدند؛ کم‌بار بودن نیز به معنی سالم بودن فرد نیست.
\newpage
\section{انتخاب قطعات و مسیر پردازش}
کانال \lr{MLII} بر اساس نام آن انتخاب شد. رکوردهای دارای ضربان‌ساز، با شناسه‌های ۱۰۲، ۱۰۴، ۱۰۷ و ۲۱۷ حذف شدند. دو رکورد ۲۰۱ و ۲۰۲ متعلق به یک فرد در نظر گرفته شدند. برای ورود هر قطعه، حداقل \fnum{۹۹٪} ضربان‌ها باید از انواع \lr{N, L, R, V} و حداقل \fnum{۹۵٪} مدت آن در ریتم‌های \lr{N, B, T, SBR} می‌بود. کامل‌بودن قطعه، متناهی‌بودن مقادیر و نبود خط صاف پیوستهٔ بیش از دو ثانیه نیز بررسی شد. در پایان، ۵۸ قطعه حذف و ۳۸ قطعه پذیرفته شد. علت‌های حذف می‌توانستند در یک قطعه هم‌زمان رخ دهند.
\begin{figure}[H]
\centering
\begin{LTR}
\begin{tikzpicture}[font=\small\persianfont,
base/.style={draw=navy!70,fill=pale,rounded corners=3pt,line width=.6pt,minimum height=10mm,align=center,inner sep=5pt},
arr/.style={-{Latex[length=2.1mm]},draw=navy,line width=.8pt}]
\node[base,text width=14.0cm] (data) at (0,0) {\rl{انتخاب ۳۸ قطعهٔ واجد شرایط از نوار قلب و حاشیه‌نویسی مرجع}};
\node[base,text width=6.35cm] (ref) at (3.85,-1.45) {\rl{مسیر مرجع}\\[-1pt]\rl{خواندن زمان ضربان‌ها از حاشیه‌نویسی}};
\node[base,text width=6.35cm,draw=amber,fill=amber!9] (det) at (-3.85,-1.45) {\rl{مسیر آشکارساز}\\[-1pt]\rl{پالایش نوار قلب و آشکارسازی قله‌های \lr{R}}};
\node[base,text width=14.0cm] (hr) at (0,-3.05) {\rl{فاصلهٔ ضربان‌ها \lr{$RR$}، تبدیل به \lr{$HR=60/RR$} و درون‌یابی در ۴ هرتز}};
\node[base,text width=14.0cm] (feat) at (0,-4.45) {\rl{تخمین بای‌اسپکتروم و استخراج هفت ویژگی \lr{HOS}}\\[-1pt]\rl{استخراج هشت ویژگی ساده و تشکیل مجموعهٔ ترکیبی}};
\node[base,text width=14.0cm] (eval) at (0,-5.95) {\rl{طبقه‌بندی با جداسازی افراد، پیش‌بینی خارج از آموزش و بازهٔ اطمینان}};
\draw[arr] (data.south) -- ++(0,-.2) -| (ref.north);
\draw[arr] (data.south) -- ++(0,-.2) -| (det.north);
\draw[arr] (ref.south) -- ++(0,-.25) -| (hr.north);
\draw[arr] (det.south) -- ++(0,-.25) -| (hr.north);
\draw[arr] (hr) -- (feat);
\draw[arr] (feat) -- (eval);
\end{tikzpicture}
\end{LTR}
\caption{روند کلی آزمایش. برچسب هر قطعه فقط از نوع ضربان‌های مرجع ساخته شد و شمارش یا درصد \lr{PVC} وارد مدل نشد.}
\label{fig:flow}
\end{figure}
در مسیر مرجع، زمان ضربان‌ها مستقیماً از حاشیه‌نویسی متخصصان خوانده شد تا اثر ویژگی‌ها بدون خطای آشکارساز بررسی شود. در مسیر دوم، نوار قلبِ کامل به ۳۲۰ هرتز تبدیل و با پایین‌گذر ۳۵ هرتز، بالاگذر \fnum{۰٫۳} هرتز، ناچ ۶۰ هرتز و اصلاح خط پایه پردازش شد. آشکارسازی روی کل رکورد انجام گرفت و سپس ضربان‌های هر قطعه انتخاب شدند. خروجی آشکارساز با کمک ضربان‌های مرجع اصلاح نشد.
\begin{figure}[H]
\centering\figwide{signals.pdf}{80mm}
\caption{نمونهٔ قطعهٔ کم‌بار و پربار: پنجره‌ای ده‌ثانیه‌ای از نوار قلب در سمت چپ و سری ضربان ۱۵ دقیقه‌ای در سمت راست. پنجرهٔ نمونهٔ پربار در اطراف نخستین ضربان زودرس بطنی انتخاب شده است. نشانگرها زمان ضربان‌های مرجع و آشکارسازی‌شده را مشخص می‌کنند.}
\label{fig:signals}
\end{figure}
\newpage
\section{تبدیل سری ضربان و استخراج ویژگی‌ها}
فاصلهٔ دو ضربان متوالی برحسب ثانیه و ضربان قلب برحسب ضربان در دقیقه محاسبه شد:
\begin{equation}
RR_i=t_i-t_{i-1},\qquad HR_i=\frac{60}{RR_i}.
\end{equation}
هر فاصله به زمان ضربان پایانی آن نسبت داده شد. سری \lr{HR} با درون‌یابی خطی روی شبکهٔ ۴ هرتز قرار گرفت؛ خارج از محدودهٔ زمان‌های موجود برون‌یابی نشد. فاصله‌های مربوط به ضربان زودرس و مکث پس از آن حفظ شدند، زیرا بخشی از پدیدهٔ مورد طبقه‌بندی بودند.

برای تخمین بای‌اسپکتروم، سری هر قطعه به بلوک‌های ۱۰۲۴ نمونه‌ای با هم‌پوشانی ۵۱۲ نمونه تقسیم شد. میانگین هر بلوک کم و پنجرهٔ \lr{Hann} اعمال شد. اگر \(X_m[k]\) تبدیل فوریهٔ بلوک \(m\) باشد، تخمین از رابطهٔ زیر به دست آمد:
\begin{equation}
\widehat B[k_1,k_2]=\frac{1}{K}\sum_{m=1}^{K}
X_m[k_1]X_m[k_2]X_m^*[k_1+k_2].
\end{equation}
میانگین‌گیری پیش از گرفتن قدرمطلق و روی حاصل‌ضرب‌های مختلط انجام شد. ناحیهٔ \(k_2\leq k_1\) و \(k_1+k_2\leq512\)، با حفظ مؤلفهٔ صفر، بررسی شد. کف دامنه برای محاسبهٔ لگاریتم‌ها \lr{$10^{-12}$} بود. هر ۳۸ قطعه در هر دو مسیر شش بلوک معتبر داشتند و از هر قطعه فقط یک بردار ویژگی ساخته شد.
\begin{table}[H]
\centering\small
\caption{هفت ویژگی بای‌اسپکتروم و مفهوم آن‌ها.}
\begin{tabularx}{\linewidth}{r>{\raggedleft\arraybackslash}X}
\toprule
ویژگی & تعریف مورد استفاده\\
\midrule
\lr{$P_1,P_2$} & آنتروپی شانونِ توزیع قدرمطلق و قدرمطلقِ به‌توان دو؛ نرمال‌شده با لگاریتم تعداد نقاط ناحیه.\\
\lr{$H_1,H_2,H_3$} & مجموع لگاریتم طبیعی دامنه در کل ناحیه، روی قطر و روی قطر با وزن شمارهٔ نقطه؛ برای دامنه کف عددی در نظر گرفته شد.\\
\lr{$f_{1m},f_{2m}$} & مختصات مرکز وزنی قدرمطلق بای‌اسپکتروم؛ برحسب شمارهٔ نقطهٔ طیفی با مبدأ یک، نه هرتز.\\
\bottomrule
\end{tabularx}
\label{tab:hos}
\end{table}
\begin{figure}[H]
\centering\figwide{bispectrum.pdf}{65mm}
\caption{لگاریتم قدرمطلق بای‌اسپکتروم برای همان دو نمونهٔ شکل \ref{fig:signals}. محورهای بسامد برحسب هرتز هستند. شکل‌ها ناحیهٔ اصلی تخمین را نمایش می‌دهند و مقیاس رنگ برای مقایسهٔ دو نمونه یکسان است.}
\label{fig:bispectrum}
\end{figure}
هشت ویژگی ساده نیز محاسبه شد: میانگین و انحراف معیار نمونه‌ای فاصله‌ها، ریشهٔ میانگین مربعات اختلاف‌های متوالی \lr{(RMSSD)}، سهم اختلاف‌های متوالی که قدرمطلق آن‌ها بیش از ۵۰ میلی‌ثانیه است، لگاریتم توان سه باند بسامدی و آنتروپی نرمال‌شدهٔ طیف توان. باندهای \lr{VLF}، \lr{LF} و \lr{HF} به‌ترتیب \fnum{۰٫۰۰۳۳} تا \fnum{۰٫۰۴}، \fnum{۰٫۰۴} تا \fnum{۰٫۱۵} و \fnum{۰٫۱۵} تا \fnum{۰٫۴} هرتز بودند. توان از سری ضربان قلب برآورد شد.

این شاخص‌ها تمام ضربان‌ها را شامل می‌شوند و نباید شاخص‌های صرفاً مبتنی بر ضربان طبیعی تلقی شوند. مجموعهٔ ترکیبی شامل هر ۱۵ ویژگی بود.
\newpage
\section{آموزش، ارزیابی و نتایج طبقه‌بندی}
چهار روش مقایسه شدند: پیش‌بینی کلاس اکثریت، طبقه‌بندی با ویژگی‌های ساده، با هفت ویژگی \lr{HOS} و با مجموعهٔ ترکیبی. سه روش اخیر از ماشین بردار پشتیبان با هستهٔ شعاعی \lr{(RBF-SVM)} و وزن متوازن کلاس‌ها استفاده کردند. مقادیر \(C\in\{1,10,100,190\}\) و \(\gamma\in\{\mathrm{scale},0.01,0.1\}\) فقط با دادهٔ آموزش بررسی شدند.

اعتبارسنجی متقاطع سه بخش بیرونی برای آزمون و سه بخش درونی برای تنظیم مدل داشت. تقسیم‌بندی در سطح فرد انجام شد؛ تمام قطعات یک فرد در یک سمت تقسیم ماندند. استانداردسازی فقط روی آموزشِ همان مرحله انجام گرفت. همهٔ روش‌ها از قطعات و تقسیم‌های یکسان استفاده کردند. بنابراین، هر قطعه در هر روش یک پیش‌بینی خارج از آموزش داشت.

معیار اصلی دقت متوازن بود؛ یعنی میانگین حساسیت کلاس پربار و ویژگی کلاس کم‌بار. بازه‌های اطمینان \fnum{۹۵٪} از ۲۰۰۰ بازنمونه‌گیری در سطح فرد محاسبه شدند. این بازه‌ها عدم‌قطعیتِ پیش‌بینی‌های خارج از آموزشِ موجود را توصیف می‌کنند و شامل تکرار کامل فرایند آموزش نیستند.
\begin{table}[H]
\centering\small
\caption{نتایج طبقه‌بندی؛ تمام شاخص‌ها برحسب درصد هستند. حساسیت و \lr{$F_1$} مربوط به کلاس پربار و ویژگی مربوط به کلاس کم‌بار است.}
\begin{tabular}{rrccccc}
\toprule
مسیر & ویژگی/روش & \shortstack{دقت\\متوازن} & بازهٔ \fnum{۹۵٪} & حساسیت & ویژگی & \lr{$F_1$}\\
\midrule
مرجع & اکثریت & \fnum{۵۰٫۰۰} & \interval{50.00}{50.00} & \fnum{۰٫۰۰} & \fnum{۱۰۰٫۰۰} & \fnum{۰٫۰۰}\\
 & ساده & \fnum{۹۵٫۸۳} & \interval{83.33}{100} & \fnum{۹۱٫۶۷} & \fnum{۱۰۰٫۰۰} & \fnum{۹۵٫۶۵}\\
 & \lr{HOS} & \fnum{۸۹٫۷۴} & \interval{78.92}{100} & \fnum{۸۳٫۳۳} & \fnum{۹۶٫۱۵} & \fnum{۸۶٫۹۶}\\
 & ترکیبی & \fnum{۹۳٫۹۱} & \interval{81.41}{100} & \fnum{۹۱٫۶۷} & \fnum{۹۶٫۱۵} & \fnum{۹۱٫۶۷}\\
\midrule
آشکارساز & اکثریت & \fnum{۵۰٫۰۰} & \interval{50.00}{50.00} & \fnum{۰٫۰۰} & \fnum{۱۰۰٫۰۰} & \fnum{۰٫۰۰}\\
 & ساده & \fnum{۸۷٫۸۲} & \interval{72.59}{98.28} & \fnum{۸۳٫۳۳} & \fnum{۹۲٫۳۱} & \fnum{۸۳٫۳۳}\\
 & \lr{HOS} & \fnum{۹۶٫۱۵} & \interval{89.47}{100} & \fnum{۱۰۰٫۰۰} & \fnum{۹۲٫۳۱} & \fnum{۹۲٫۳۱}\\
 & ترکیبی & \fnum{۹۶٫۱۵} & \interval{89.47}{100} & \fnum{۱۰۰٫۰۰} & \fnum{۹۲٫۳۱} & \fnum{۹۲٫۳۱}\\
\bottomrule
\end{tabular}
\label{tab:results}
\end{table}
\begin{figure}[H]
\centering\figwide{performance.pdf}{79mm}
\caption{دقت متوازنِ چهار روش در دو مسیر زمانی. خط‌چین سطح \fnum{۵۰٪} و میله‌های خطا بازهٔ اطمینان \fnum{۹۵٪} را نشان می‌دهند. مقایسه بر پایهٔ همان افراد و همان قطعات انجام شده است.}
\label{fig:performance}
\end{figure}
در آزمایش اصلیِ مسیر مرجع، مدل ساده بالاترین دقت متوازن را داشت. اختلاف \lr{HOS} منهای مدل ساده برابر با \fnum{۶٫۰۹} واحد درصد کمتر، با بازهٔ اطمینان \interval{-19.86}{+11.54} واحد درصد بود. در مسیر آشکارساز، این اختلاف \fnum{۸٫۳۳} واحد درصد بیشتر و بازهٔ آن \interval{0.00}{+21.45} بود. این نتایج به‌تنهایی برتری پایدار یک مجموعهٔ ویژگی را تأیید نمی‌کنند.
\newpage
\section{کیفیت آشکارسازی ضربان و الگوی خطاها}
ضربان‌های آشکارسازی‌شده با مرجع به‌صورت یک‌به‌یک و با تلورانس ۱۵۰ میلی‌ثانیه تطبیق داده شدند. در مجموع قطعات پذیرفته‌شده، ۴۰۳۸۸ ضربان تطبیق یافت، ۲۳۱ تشخیص اضافی و ۵۲۲ ضربان ازدست‌رفته وجود داشت. خطای زمانی تنها برای ضربان‌های تطبیق‌یافته محاسبه شد.
\begin{table}[H]
\centering\small
\caption{کیفیت تجمیعی آشکارسازی ضربان در قطعات پذیرفته‌شده.}
\begin{tabular}{cccc}
\toprule
حساسیت & ارزش اخباری مثبت & \lr{$F_1$} & میانگین خطای مطلق زمانی\\
\midrule
\fnum{۹۸٫۷۲٪} & \fnum{۹۹٫۴۳٪} & \fnum{۹۹٫۰۸٪} & \fnum{۳٫۱۸} میلی‌ثانیه\\
\bottomrule
\end{tabular}
\end{table}
این میانگین‌ها کیفیت یکسان تمام قطعات را نشان نمی‌دهند. برای نمونه، در ۱۵ دقیقهٔ اول رکورد ۲۲۸، حساسیت آشکارسازی فقط \fnum{۶۵٫۱۶٪} بود. در نیمهٔ دوم رکورد ۱۰۸ نیز هم ضربان ازدست‌رفته و هم تشخیص اضافی دیده شد. پس از تبدیل این زمان‌ها به سری ضربان، این خطاها می‌توانند روی ویژگی‌های ساده و بای‌اسپکتروم اثر بگذارند.
\begin{table}[H]
\centering\small
\caption{سه قطعه با کمترین حساسیت آشکارسازی؛ بازهٔ زمانی نسبت به آغاز رکورد است.}
\begin{tabular}{rccccc}
\toprule
رکورد و بازه & حساسیت & ارزش اخباری مثبت & تطبیق‌یافته & اضافی & ازدست‌رفته\\
\midrule
۲۲۸، دقیقهٔ ۰ تا ۱۵ & \fnum{۶۵٫۱۶٪} & \fnum{۹۸٫۴۰٪} & ۶۷۷ & ۱۱ & ۳۶۲\\
۱۰۸، دقیقهٔ ۱۵ تا ۳۰ & \fnum{۸۹٫۶۲٪} & \fnum{۸۳٫۳۳٪} & ۸۲۰ & ۱۶۴ & ۹۵\\
۲۲۸، دقیقهٔ ۱۵ تا ۳۰ & \fnum{۹۶٫۵۳٪} & \fnum{۹۹٫۹۰٪} & ۹۷۳ & ۱ & ۳۵\\
\bottomrule
\end{tabular}
\end{table}
\begin{figure}[H]
\centering\figwide{confusion.pdf}{71mm}
\caption{ماتریس درهم‌ریختگی مدل \lr{HOS} در دو مسیر. هر خانه تعداد قطعات را نشان می‌دهد؛ سطرها کلاس مرجع و ستون‌ها پیش‌بینی مدل هستند. هیچ قطعه‌ای در مرحلهٔ استخراج ویژگی شکست نخورد.}
\label{fig:confusion}
\end{figure}
در مسیر مرجع، مدل \lr{HOS} از ۲۶ قطعهٔ کم‌بار، ۲۵ قطعه و از ۱۲ قطعهٔ پربار، ۱۰ قطعه را درست تشخیص داد. در مسیر آشکارساز، هر ۱۲ قطعهٔ پربار درست پیش‌بینی شدند، اما دو قطعهٔ کم‌بار به‌اشتباه پربار شناخته شدند. در هر دو مسیر، استخراج ویژگی برای تمام قطعات پذیرفته‌شده موفق بود؛ با این حال، موفقیت استخراج به معنی دقیق بودن همهٔ زمان‌های آشکارسازی نیست.

بالاتر بودن دقت طبقه‌بندیِ \lr{HOS} در مسیر آشکارساز، نشان نمی‌دهد که زمان ضربان‌های آن از مرجع دقیق‌تر بوده است. اختلاف دقت متوازن این دو مسیر \fnum{۶٫۴۱} واحد درصد بود و بازهٔ اطمینان جفت‌شدهٔ آن، از \fnum{۳٫۸۵} واحد درصد کمتر تا \fnum{۱۶٫۶۷} واحد درصد بیشتر، صفر را دربر می‌گرفت.
\newpage
\section{بحث}
نتیجهٔ اصلی این آزمایش آن است که در مسیر مرجع، ویژگی‌های سادهٔ سری ضربان دست‌کم در برآورد نقطه‌ای از \lr{HOS} بهتر عمل کردند. این موضوع با ماهیت هدف آزمایش سازگار است: بار ضربان زودرس می‌تواند الگوی فاصله‌های کوتاه و بلند را تغییر دهد و شاخص‌هایی مانند پراکندگی فاصله‌ها و اختلاف‌های متوالی، بخشی از این تغییر را مستقیماً بازتاب دهند. افزودن هفت ویژگی \lr{HOS} نیز در مسیر مرجع، عملکرد مجموعهٔ ساده را افزایش نداد. این توضیح، تفسیر روش‌شناختیِ نتیجه است و سازوکار فیزیولوژیک مستقلی را اثبات نمی‌کند.
\begin{figure}[H]
\centering\figwide{features.pdf}{80mm}
\caption{توزیع هفت ویژگی \lr{HOS} در مسیر مرجع. هر نقطه یک قطعهٔ ۱۵ دقیقه‌ای است؛ نقاط متعلق به یک فرد مستقل از هم فرض نمی‌شوند. این نمایش توصیفی است و برای انتخاب ویژگی یا تغییر قواعد آزمایش استفاده نشد.}
\label{fig:features}
\end{figure}
مسیر آشکارساز نتیجهٔ متفاوتی داشت و \lr{HOS} و مجموعهٔ ترکیبی در آن به دقت متوازن \fnum{۹۶٫۱۵٪} رسیدند. با این حال، خطای زمان‌یابی ضربان‌ها نیز ساختار سری ضربان را تغییر می‌دهد. در این مطالعه نمی‌توان سهم اطلاعات مربوط به ریتم را از اثر خطای آشکارساز به‌طور قطعی جدا کرد. بنابراین، نتیجهٔ این مسیر باید همراه با کیفیت آشکارسازی هر قطعه و مقایسهٔ جفت‌شده با مسیر مرجع تفسیر شود.

محدودیت اصلی، وجود تنها ۲۲ فرد و به‌ویژه هفت فرد در گروه پربار است. تعداد قطعات بیشتر از تعداد افراد بود، اما جداسازی و بازنمونه‌گیری در سطح فرد انجام شد تا قطعات یک فرد به‌عنوان افراد مستقل شمرده نشوند. ارزیابی روی پایگاه مستقلی انجام نشد و بازه‌های اطمینان نیز مشروط به پیش‌بینی‌های همین اعتبارسنجی هستند. علاوه بر این، کنار گذاشتن بارهای بین \fnum{۱٪} و \fnum{۵٪} باعث می‌شود نتیجه مستقیماً به قطعات مرزی تعمیم‌پذیر نباشد. این مطالعه دربارهٔ دو سطح بار \lr{PVC} است و نتیجهٔ آن را نمی‌توان به تشخیص پنج گروه بیماریِ مطالعهٔ مبنا تعمیم داد.
\section{جمع‌بندی}
فرایند استخراج ویژگی بای‌اسپکتروم و طبقه‌بندی بار \lr{PVC} روی دادهٔ عمومی با موفقیت اجرا شد. هر دو مسیر برای تمام ۳۸ قطعه ویژگی تولید کردند و ارزیابی با جداسازی افراد انجام شد. در مسیر مرجع، ویژگی‌های ساده با دقت متوازن \fnum{۹۵٫۸۳٪} نسبت به \lr{HOS} با \fnum{۸۹٫۷۴٪} دقت متوازن بالاتری داشتند؛ نتیجهٔ \fnum{۹۶٫۱۵٪} در مسیر آشکارساز نیز با وجود خطاهای موضعی و عدم‌قطعیت باید تفسیر شود. حاصل این کار، یک نتیجهٔ اکتشافی دربارهٔ امکان استفاده از \lr{HOS} است؛ شواهد حاضر افزایش پایدار عملکرد نسبت به ویژگی‌های ساده را تأیید نمی‌کنند.
\begingroup\footnotesize
\begin{thebibliography}{9}
\begin{latin}
\bibitem{chua} Chua, K. C., Chandran, V., Acharya, R. U., \& Lim, C. M. (2009).
Cardiac Health Diagnosis Using Higher Order Spectra and Support Vector Machine.
\textit{The Open Medical Informatics Journal}, 3, 1--8.
\href{https://doi.org/10.2174/1874431100903010001}{doi:10.2174/1874431100903010001}.
\bibitem{mitdb} Moody, G., \& Mark, R. (2005).
\textit{MIT-BIH Arrhythmia Database} (version 1.0.0). PhysioNet.
\href{https://doi.org/10.13026/C2F305}{doi:10.13026/C2F305}.
\end{latin}
\end{thebibliography}
\endgroup
\end{document}
